\subsection*{One local zero of the complete ground transform at $\lambda=4$}

The finite-compression semantic lock does not itself locate complete-ground
zeros. The following transfer uses the complete even-sector spectral
separation and the energy of an exact archived trial. No new window or
finite eigenvector is computed.

\begin{lemma}[Energy control of a ground projection]
\label{lem:v142-energy-projection}
Let $W\succeq0$ be self-adjoint on a real Hilbert space of functions in
$L^2([-a,a])$, with a simple normalized ground $u$. Suppose an invariant
closed subspace contains $u$ and its remaining spectrum lies in
$[b,\infty)$, where $b>0$. For a real unit trial $f$ in that subspace and
in the form domain, put $\rho=q_W[f]<b$, $h=P_u f$, and
$\eta=\sqrt{\rho/b}$. Then $h\ne0$ and
\[
 \norm{f-h}_2\le\eta.
\]
For the nonunitary Fourier transform and every real $t$,
\begin{equation}\label{eq:v142-transform-budget}
 |\widehat f(t)-\widehat h(t)|\le\sqrt{2a}\,\eta,
 \qquad
 |\widehat f'(t)-\widehat h'(t)|\le\sqrt{2a^3/3}\,\eta.
\end{equation}
If both functions are even, opposite endpoint signs for $\widehat h$
and a strictly positive derivative throughout an interval prove a
unique simple zero there, also a zero of $\widehat u$.
\end{lemma}
\begin{proof}
The spectral theorem and nonnegativity give
$q_W[f]\ge b\norm{(I-P_u)f}_2^2$. Since $\eta<1$, the orthogonal
projection $h$ is nonzero. Cauchy--Schwarz applied to $e^{-itx}$ and
$xe^{-itx}$ gives \eqref{eq:v142-transform-budget}. Compact support
justifies differentiation and makes the transforms entire. Reality
and evenness make them real on the real axis. The intermediate value
theorem and strict monotonicity give the asserted unique zero, whose
nonzero derivative makes it simple. Finally $h$ is a nonzero scalar
multiple of $u$. This comparison uses $h=P_u f$, not a separately
renormalized approximation to $u$; no normalization error is omitted.
\end{proof}

\begin{proposition}[A complete-ground local zero near the first zeta ordinate]
\label{prop:v142-ground-zero}
Let $\gamma_1$ be the first positive ordinate of a Riemann zeta zero,
and let $\xi_4$ be the complete ground of
Corollary~\ref{cor:v140-real-zeros}. Then $\widehat\xi_4$ has exactly
one zero in the real interval
\[
 I=(\gamma_1-1/10,\gamma_1+1/10),
 \qquad \gamma_1=14.13472514173469\ldots,
\]
and that zero is simple. This is a local statement: it does not order
all earlier zeros, determine the sign of the discrepancy, identify
the zero with $\gamma_1$, or establish convergence to Riemann's $\Xi$.
\end{proposition}
\begin{proof}[Complete separation and two-precision interval gates]
Use the complete even-sector data of
Proposition~\ref{prop:v140-ground4} without changing the physical window,
the head cut $16$, the exact trial support $4096$, or the archived
complete residual bounds. Apply Lemma~\ref{lem:v140-shifted-schur} with
\[
 \kappa=62629/100000,\quad \delta=1.7940\cdot10^{-8},
 \quad b=10^{-67}.
\]
At both $1024$ and $1280$ bits, the $17$-dimensional lower matrix
$L_b$ has signed LDL pivots
\[
 \underbrace{+\,,\ldots,\,+}_{15},\ -\,,\ +.
\]
Its positive subspace has dimension $16$. The exact even trial below
$b$ supplies a strictly negative direction of the complete shifted
form, while $T-bI\succ0$. Square completion and compact resolvent,
as in Proposition~\ref{prop:v140-ground4}, show that precisely one
complete even eigenvalue is below $b$ and none equals $b$.
Thus every even vector orthogonal to the ground has spectral support
strictly above $10^{-67}$. This is an even-sector separator, not an
improved bound for the first excited eigenvalue across both parities.

Take the same real trial column $G_{\rm ev}e_{16}$, divided by its
ordinary norm, and center it on $[-\log4,\log4]$. Its complete
Rayleigh quotient is
\[
 2.45361754930\cdot10^{-75}<\rho<2.45361754931\cdot10^{-75}.
\]
With $L=2\log4$, Lemma~\ref{lem:v142-energy-projection} gives the
following outward bounds for $h=P_{\xi_4}f$:
\begin{align*}
 \norm{f-h}_2&<1.566403\cdot10^{-4},\\
 E_0=\sqrt L\sqrt{\rho/b}&<0.000260824,\\
 E_1=\sqrt{L^3/12}\sqrt{\rho/b}&<0.000208757.
\end{align*}
All $4097$ exact trial coefficients are retained in
\eqref{eq:v141-transform}, with its correct centering convention.
Interval evaluation at Arb enclosures of $\gamma_1\pm1/10$ gives
\[
 \widehat f(\gamma_1-1/10)+E_0<-0.00006731,
 \qquad
 \widehat f(\gamma_1+1/10)-E_0>0.00003349.
\]
The differentiated formula, evaluated on eight covering subintervals
of length $1/40$, gives
\[
 \widehat f'(t)-E_1>0.0016\quad(t\in\overline I).
\]
These three strict inequalities transfer to $\widehat h$ and prove
the result. Multiplication by the unitary Fourier factor or by the
nonzero ground-projection scalar preserves the zero and its simplicity.

Both precision runs reconstruct the exact dyadic trial and ordinary
Grams. The inherited even residual enclosures at $768/896$ bits
include every row beyond $J=65536$; their small complete Schur gate
is replayed, but their full assemblies are not recomputed.
A second implementation checks the shifted negative index with
principal determinants and the transform with termwise centered
cosine integrals, including derivative signs on $32$ subintervals.
Saved decimal intervals are reparsed to check that serialization
preserves all gates. Scripts, input hashes and reports are in
\texttt{evidence/v142/}. This implementation cross-check is not an
external audit.
\end{proof}

\paragraph{What had to change, and what remains open.}
Using the old global separator $10^{-73}$ in the same energy comparison
would give $E_0\simeq0.260823$, too large to certify either endpoint
sign. Restricting the comparison to the invariant even sector permits
$10^{-67}$ and reduces this error bound by a factor of $1000$.
The obstruction was the coarse separation bound, not a missing zero
of the complete ground transform. The resulting $0.1$ enclosure
still does not transfer the minute CCM finite-compression discrepancies.
A sharper energy or residual-to-transform estimate would be needed
for that task. Cofinal convergence, G2 and RH remain open.
